\chapter{Chamada dos arquivos dentro do código principal}
O arquivo principal a ser compilado é {\ttfamily PPGMC.tex}. Nele são chamados os arquivos/elementos da sua dissertação (ou tese). Faça as modificações após o comando \verb#\begin{document}#. Depois, abra cada um dos documentos que estão sendo chamados pelo documento mestre e modifique-os para obter o seu documento final. Todos os arquivos dos elementos pré-textuais, textuais e pós-textuais devem ser editados pelos autores, segundo suas necessidades.

\section{Arquivos, Pacotes e configurações}
Na pasta "pacotes" está o arquivo {\ttfamily ppgmc-uesc} que contém toda a programação da classe. \textbf{Esses arquivos não devem ser apagados ou alterados}.

O comando \verb#\pretextual# indica o início da chamada dos elementos obrigatórios e opcionais. Os comandos \verb#\imprimircapa# e \verb#\imprimirfolhaderosto{}# imprimem a capa e a folha de rosto, respectivamente, com as informações contidas no arquivo {\ttfamily 1-Inf.Capa-FolhaRosto} na pasta {\ttfamily 01-elementos-pre-textuais}. Tanto a capa quanto a folha de rosto são obrigatórias segundo as normas da ABNT e da UESC.

\textit{No arquivo \textit{{\ttfamily 1-Inf.Capa-FolhaRosto}} referido é que você deve colocar as informações principais do seu trabalho, tais como: título, autor, orientador, coorientador (caso não tenha coorientador, comente a linha), instituição, tipo de trabalho (dissertação ou tese) etc. Tais informações aparecerão automaticamente na capa, folha de rosto e folha de aprovação.}

Este template leva em consideração que o autor possui coorientador. Caso não possua, basta comentar a linha 43 do arquivo \textit{{\ttfamily FolhaAprovacao}}, também na pasta {\ttfamily 01-elementos-pre-textuais}.

Na pasta {\ttfamily 01-elementos-pre-textuais} estão exemplos de arquivos pré-textuais, chamados no arquivo principal, a saber:
\begin{alineas}
\item \textbf{FolhaAprovacao} -- elemento obrigatório
\item \textbf{dedicatoria} -- elemento opcional
\item \textbf{agradecimentos} -- elemento opcional
\item \textbf{epigrafe} -- elemento opcional
\item \textbf{resumoPt} -- elemento obrigatório
\item \textbf{resumoEn} -- elemento obrigatório
\item \textbf{listaFiguras} -- elemento opcional
\item \textbf{listaTabelas} -- elemento opcional
\item \textbf{listaQuadros} -- elemento opcional
\item \textbf{listaAlgoritmos} -- elemento opcional
\item \textbf{listaSiglas} -- elemento opcional
\item \textbf{listaSimbolos} -- elemento opcional
\item \textbf{sumário} -- elemento obrigatório
\end{alineas}

Os arquivos {\ttfamily listaFiguras,\! listaTabelas,\! listaQuadros e\! listaAlgoritmos}, não precisam ser modificados. Caso algum desses elementos não seja necessário no seu trabalho, basta comentar a linha correspondente no arquivo principal. Já os arquivos {\ttfamily listaSiglas} e {\ttfamily listaSimbolos} devem ser editados caso façam parte do seu trabalho. Caso contrário, basta comentar a linha correspondente no arquivo principal.

A \textit{ficha catalográfica} é um elemento obrigatório fornecido pela biblioteca e deve ser colocado no verso da folha de rosto. A linha correspondente no arquivo principal está comentada. Após receber a ficha catalográfica da biblioteca, faça upload do arquivo renomeado para \textit{{\ttfamily FichaCatalografica}} e descomente a linha correspondente no arquivo principal.

O comando \verb#\textual# marca o início dos elementos textuais, que incluem a introdução, os capítulos e a conclusão. Já o comando \verb#\postextual# indica o início dos elementos pós-textuais, que incluem as referências, o glossário, os apêndices e os índices. As referências devem ser armazenadas em um arquivo com extensão \verb#.bib#. O arquivo {\ttfamily refbase.bib}, que contém uma quantidade de referências como exemplo, está localizado na pasta \verb#03-elementos-pos-textuais#. Você pode substituir as referências desse arquivo pelas suas, sem precisar criar um novo arquivo. Mais detalhes sobre isso estão no Capítulo \ref{chap:conclusao}.
\section{Coleta de dados}

Inserir seu texto aqui...

